\section{System Architecture}
\label{sec:architecture}

The system has three core services and one control process, all on a single host. They never call each other in the
default mode: they exchange messages through Hazelcast and keep their data in the datalake and in MongoDB.
Figure~\ref{fig:architecture} shows the whole deployment: the front-end boundary, the internal interactions, the
replicas and threads of each process, and the shared storage.

\begin{figure}[ht]
\centering
\begin{tikzpicture}[
  font=\footnotesize,
  >={Stealth[length=4pt]},
  box/.style={draw=inksecondary, rounded corners=2pt, fill=white, align=center, inner sep=3pt},
  service/.style={box, minimum width=3.05cm, minimum height=1.25cm},
  queue/.style={draw=modehazelcast, rounded corners=1pt, fill=white, align=center, minimum width=2.25cm, minimum height=0.55cm, inner sep=2pt},
  topic/.style={draw=modehazelcast, ellipse, fill=white, align=center, inner sep=1pt, minimum width=2.4cm, minimum height=0.6cm},
  store/.style={draw=inksecondary, fill=white, rounded corners=1pt, align=left, inner sep=3pt, anchor=west},
  frontend/.style={->, draw=modehazelcast!80!black, line width=1.1pt},
  operation/.style={->, draw=inksecondary, line width=0.8pt},
  message/.style={->, draw=modehazelcast, line width=0.8pt, dashed},
  access/.style={draw=inkprimary!70, line width=0.5pt},
  lbl/.style={font=\scriptsize, text=inksecondary, inner sep=1pt},
]
\fill[hostfill, rounded corners=4pt] (0,0) rectangle (15.6,9.55);
\draw[inkmuted, rounded corners=4pt] (0,0) rectangle (15.6,9.55);
\node[anchor=north west, font=\footnotesize\bfseries, text=inksecondary] at (0.15,9.5) {One host (Docker Compose)};

\node[box, minimum width=13.6cm, minimum height=0.6cm] (nginx) at (7.8,8.55)
  {\textbf{nginx} $\times$1 \quad 8080 front-end \quad 8081 operations \quad 8082 internal routes (mode \mode{orchestrator} only)};

\node[box] (admin) at (5.65,10.75) {Scripts and admin};
\node[box] (fe) at (12.4,10.75) {Front-end (browser)};
\draw[operation] (admin.south) -- node[lbl, right, pos=0.3] {port 8081} (admin.south |- nginx.north);
\draw[frontend] (fe.south) -- node[lbl, left, pos=0.3] {port 8080: the front-end API boundary} (fe.south |- nginx.north);

\node[service] (control) at (1.95,6.55) {\textbf{control} $\times$1 \quad :7004\\reconciliation thread\\HTTP pool};
\node[service] (crawler) at (5.65,6.55) {\textbf{crawler} $\times$N \quad :7001\\8 download threads\\HTTP pool};
\node[service] (indexer) at (9.35,6.55) {\textbf{indexer} $\times$N \quad :7002\\1 batch thread\\HTTP pool};
\node[service] (query) at (13.05,6.55) {\textbf{query} $\times$N \quad :7003\\1 snapshot thread\\HTTP pool, read only};

\draw[operation] (control.north |- nginx.south) -- node[lbl, right] {/control} (control.north);
\draw[operation] (crawler.north |- nginx.south) -- node[lbl, right] {/ingest} (crawler.north);
\draw[operation] (indexer.north |- nginx.south) -- node[lbl, right] {/index} (indexer.north);
\draw[frontend] (query.north |- nginx.south) -- node[lbl, left, align=right] {/search, /books,\\/search/filters} (query.north);

\begin{scope}[on background layer]
\fill[white, rounded corners=3pt] (2.45,3.7) rectangle (12.55,5.15);
\draw[modehazelcast, rounded corners=3pt] (2.45,3.7) rectangle (12.55,5.15);
\end{scope}
\node[anchor=south east, font=\scriptsize\bfseries, text=modehazelcast] at (12.5,3.74) {Hazelcast $\times$1 \quad :5701};
\node[queue] (b2d) at (3.8,4.55) {books-to-download};
\node[queue] (b2i) at (7.5,4.55) {books-to-index};
\node[topic] (topic) at (11.2,4.55) {published-batches};

\draw[message] ($(control.south east)+(-0.25,0)$) -- node[lbl, right, pos=0.4] {seed} (b2d.north west);
\draw[message] (b2d.north east) -- ($(crawler.south west)+(0.25,0)$);
\draw[message] ($(crawler.south east)+(-0.25,0)$) -- node[lbl, right, pos=0.4] {BookTask} (b2i.north west);
\draw[message] (b2i.north east) -- ($(indexer.south west)+(0.25,0)$);
\draw[message] ($(indexer.south east)+(-0.25,0)$) -- node[lbl, right, pos=0.45, align=left] {Batch\\Published} (topic.north west);
\draw[message] (topic.north east) -- ($(query.south west)+(0.25,0)$);

\foreach \x in {2.25, 5.95, 9.65} {
  \draw[access] (\x,5.92) -- (\x,1.25);
  \fill[inkprimary!70] (\x,1.25) circle (1.3pt);
}
\draw[access] (13.05,5.92) -- (13.05,1.25);
\fill[inkprimary!70] (13.05,1.25) circle (1.3pt);
\draw[access, ->] (1.65,5.92) -- node[lbl, left, pos=0.84] {.tmp} (1.65,2.72);
\draw[access, ->] (5.35,5.92) -- node[lbl, left, pos=0.84] {write} (5.35,2.72);
\draw[access, <-] (9.05,5.92) -- node[lbl, left, pos=0.84] {read} (9.05,2.72);
\node[store, fill=white, minimum width=10.15cm, minimum height=0.55cm] (datalake) at (0.55,2.45)
  {\textbf{datalake volume}: header and body of each book, Stage~1 \code{batch} layout};
\node[store, minimum width=14.5cm, minimum height=0.8cm] (mongo) at (0.55,0.85)
  {\textbf{MongoDB} $\times$1 \quad :27017 \quad source of truth: \code{book_states}, \code{index_state},\\
  \hphantom{\textbf{MongoDB} $\times$1 \quad :27017 \quad }\code{metadata_g<n>}, \code{postings_g<n>}, \code{batches_g<n>}};

\end{tikzpicture}
\caption{Deployment on the single Stage~2 host, in the default mode \mode{hazelcast}. Thick blue arrows are front-end
requests through nginx port 8080, gray arrows are operational requests through port 8081, dashed arrows are Hazelcast
messages and thin vertical lines are storage access; every service reads and writes MongoDB. ``$\times$N'' marks a
service that scales with \texttt{docker compose up -{}-scale}, and each box lists the threads of one replica. Not
drawn: \code{crawler} also re-enqueues retried downloads in \code{books-to-download}, and \code{control} re-enqueues
lost messages and reindexing in both queues (\spec{4.2}).}
\label{fig:architecture}
\end{figure}

\subsection{Services and shared storage}

\begin{itemize}
\item \textbf{\code{crawler} (Ingestion Service)} takes \code{BookTask} messages from \code{books-to-download}. It claims
  each book in MongoDB, downloads it from Gutenberg or from the local mirror and stores its header and body in the
  datalake with the Stage~1 \code{batch} layout. It marks the book \code{DOWNLOADED} and then enqueues it in
  \code{books-to-index}.
\item \textbf{\code{indexer} (Indexing Service)} groups \code{books-to-index} into batches of up to 100 books or 30\,s.
  It extracts the metadata of each book and counts its terms. It writes metadata and postings into the collections of
  the write generation and publishes the batch with one conditional write. Then it announces \code{BatchPublished}.
\item \textbf{\code{query} (Search Service)} keeps an immutable snapshot of the published batches of the current
  generation. It serves keyword search with author, language and year filters, TF-IDF ranking, book details and filter
  values.
\item \textbf{\code{control}} seeds the candidate books at startup and reconciles every 30\,s. It releases expired
  leases, abandons expired batches, requeues lost messages, drives index rebuilds and deletes orphan \code{.tmp} files.
  In the default mode it never drives a book step by step, so it is not an orchestrator.
\item \textbf{MongoDB} is the only source of truth. It holds \code{book_states}, one document per book, and
  \code{index_state}, which points to the generation \code{query} serves. Each generation of the index has three
  collections, \code{metadata_g<n>}, \code{postings_g<n>} and \code{batches_g<n>}.
\item \textbf{Hazelcast} carries two work queues and one topic. It is not a source of truth: a lost or repeated message
  never corrupts state, because the claim in MongoDB decides who owns each book.
\end{itemize}
